\section{Discussion and limitations}

\subsection{The formation advantage is a reach advantage}

$\Delta_\text{formation}$ is not a constant. It tracks how much of what a task needs
sits outside the roster, and it collapses when reach stops being scarce.

At a low external fraction the gap is a third of its high-$E$ value ($+0.112$ against
$+0.352$ cold), and on the rework beat---where the task must return to a counterpart it
has already used---the point estimate reverses to $-0.037$: the bounded network is ahead.
The formation effect halves on rework: $0.248$ at the cold beat against $0.122$ at the
rework beat, with non-overlapping intervals.

We report the reversal as suggestive rather than established: it is a subcell point
estimate and we do not attach an interval to it. But the pattern across the table is
consistent and it bounds the headline. The open network wins because it can address
holders the bounded one cannot; where the bounded roster already contains the holders,
and where the work requires going back to them, that advantage is largely spent.

\subsection{A bounded roster as a wall, and what happens when it is not}

A bounded network here addresses twenty holders and no others, which is a truncated
directory rather than a private network: information reaches a closed group through the
people in it, one hop further out. We let a roster member forward one question to a
contact of its own, minted under its own authority, and re-ran. \textbf{The gap did not
closes about half}: the hop is worth $+0.140$ (CI $[+0.073, +0.204]$) over 28 episodes
per arm, replicating independently within each of the three edge-cost strata, and it cuts
fabrication from $0.223$ to $0.057$ per task. A first estimate from a 16-episode cell put
it at $+0.042$ with an interval spanning zero; \S\ref{app:relay2} gives the current
numbers and \S\ref{app:relay} the design and the harness failure it exposed.

\subsection{Forming an edge is free here, and that is the main limitation}

Whether a link is worth its cost is the central question of strategic network formation
\citep{jackson1996}, where agents weigh the benefit of a connection against the price of
maintaining it. \textbf{We set that price to zero.} In this design, addressing a stranger
costs exactly what addressing a teammate costs.
The requester searches a directory and sends a message; there is no negotiation, no
context transfer, no format alignment, no establishing what a counterpart can actually
do.

This removes the economic reason relationships exist. An open arm with a hundred
addressable holders at zero acquisition cost is closer to an oracle than to a network,
and the formation effect we measure should be read as an upper bound: it is what open
discovery is worth \emph{when discovery is free}. A per-new-edge cost---an onboarding
exchange, a latency penalty, a token charge on first contact---is the single variable
most likely to move the ranking. We have since run it, and it does not: charging one
and then two handshake turns per contact outside the roster---578 handshakes fired---leaves
the open arm ahead by $+0.219$, $+0.263$ and $+0.244$ across the three levels while its
tokens rise from 53k to 74k. The arm converts friction into cost rather than into failure,
which is the same shape as the directory-size null (\S\ref{app:retracted}). A cap on
\emph{total} contacts, which cannot be paid off by spending more, is the manipulation that
remains untried.

\subsection{What the attribute null does and does not bound}

The null covers assertions in a prompt, the weakest form of relational encoding. It says
nothing about trust wired into routing, memory, retrieval or permission policy, where
reputation lives in deployed systems and where it changes what an agent is able to do
rather than what it is told.

It is narrower still, in a way we can name: each seeded property had \emph{no available
action} behind it. Trust and origin were asserted uniformly over every card, ranking no
counterpart above another and so unable to change who gets asked first --- deployed
reputation is differential, and that is the whole of its value. Accountability went to
responders whose knowledge is a fixed list, leaving them nothing to do differently.

A fourth manipulation removes the first defect: a per-card line reporting how often that
card's past answers held up, accurate by construction, against a placebo assigning the
same marks at random. The marks are perfect --- 100\% accurate against 9\% --- and $F_1$
still does not move ($-0.017$, CI $[-0.174, 0.143]$). The routing measurement says why:
under the accurate marks the requester contacts flagged cards at \emph{six times} the base
rate rather than avoiding them, because the flagged cards hold the planted falsehoods and
in this task those are the same cards holding the true facts on their subject. Acting on
the flag meant not asking the only source.

So the claim is not that relational metadata does not work, but that it moves behavior
only where it discriminates between options an agent can choose between --- and three
signals that did not, one uniform, one aimed at an agent without discretion, one flagging
the sole source, moved nothing. Making it actionable needs two holders per fact, one
reliable and one not: a change to the task, not to a configuration.

Four further limits, stated plainly: the cold-phase crossing follows from the fact
allocation and is not a discovery, so only its location and cost are measured;
forced-value deliverables cannot capture whether an artifact is well argued, which is the
price of removing the judge; the kernel is a reference implementation rather than a
deployment; and one model family carries every sweep.
